\chapter{Curated EasyCrypt Proof-Object Case Studies}
\label{app:curated-easycrypt-case-studies}

This appendix gives close readings of representative EasyCrypt objects.  The
archived exhaustive catalog remains in \ec{dissertation/etc/}; the purpose here
is different.  Each selected code fragment is a proof object with a mathematical
meaning in the HAETAE security argument.

\section{The generic signature interface}
\label{app:case-sig-rom-interface}

\begin{lstlisting}[language=EasyCrypt,caption={Signature objects and ROM-query objects in \ec{Sig\_ROM.ec}.},label={lst:sig-rom-types}]
type pkey.
type skey.
type message.
type context.
type signature.
type ro_query.
type ro_output.

type query = message * context.
type signed_query = query * signature.
\end{lstlisting}

Mathematically, this declaration introduces carrier sets
\[
  \mathcal{PK},\mathcal{SK},\mathcal{M},\mathcal{C},\Sigma,
  \mathcal{Q}_{RO},\mathcal{Y}_{RO}.
\]
The type \ec{query} is the product \(\mathcal{M}\times\mathcal{C}\).  The type
\ec{signed\_query} is
\((\mathcal{M}\times\mathcal{C})\times\Sigma\).  The product
\(\mathcal{M}\times\mathcal{C}\) is the substrate for EUF-CMA freshness: a final
forgery must verify and its message/context pair must not have been previously
signed.  The signed-query product is used for the stronger freshness predicate in
the strong-EUF-CMA style game, where the exact message/context/signature triple
must not have been returned before.

\begin{lstlisting}[language=EasyCrypt,caption={Scheme and adversary interfaces in \ec{Sig\_ROM.ec}.},label={lst:sig-rom-module-types}]
module type POracle = {
  proc get(q : ro_query) : ro_output
}.

module type Scheme (H : POracle) = {
  proc kg() : pkey * skey
  proc sign(sk : skey, m : message, ctx : context) : signature
  proc verify(pk : pkey, m : message, ctx : context,
              sig : signature) : bool
}.

module type SignOracle = {
  proc sign(m : message, ctx : context) : signature
}.

module type Adversary (H : POracle, O : SignOracle) = {
  proc forge(pk : pkey) : message * context * signature
}.
\end{lstlisting}

The mathematical reading is a family of kernels indexed by an oracle
implementation \(H\).  A concrete signature scheme is a functor from random
oracles to probabilistic procedures.  A chosen-message adversary is likewise a
module parameterized by the hash oracle and the signing oracle.

\section{The HAETAE scheme wrapper}
\label{app:case-haetae-scheme-wrapper}

The signing procedure in \ec{HAETAE\_Scheme.ec} makes the proof-relevant oracle
sites explicit.

\begin{lstlisting}[language=EasyCrypt,caption={Signing wrapper in \ec{HAETAE\_Scheme.ec}.},label={lst:haetae-scheme-sign-wrapper}]
coins <$ drandom_coins;
ro_y <@ H.get(sampler_expand_query coins);
coins <- ro_signing_coins ro_y;
pk <- public_key_of_secret haetae_mode sk;
ro_y <@ H.get(message_hash_query pk ctx m);
mu <- ro_message_hash ro_y;
highbits <- commitment_highbits haetae_mode sk m ctx coins;
lowbits <- commitment_lowbits haetae_mode sk m ctx coins;
ro_y <@ H.get(challenge_hash_query haetae_mode highbits lowbits mu);
sig <- sign_internal haetae_mode sk m ctx coins;
return sig;
\end{lstlisting}

This code denotes the kernel
\[
  \mathsf{Sign}^{H}(sk,m,ctx):1\to\mathcal{D}(\Sigma).
\]
The first ROM query is the sampler-expansion site, the second is the message
hash site, and the third is the challenge hash site.  The proof depends on these
sites being disjoint typed constructors; otherwise ROM programming for the
sampler could interfere with ROM programming for the challenge.

The snippet also exposes a model/spec issue.  Figure 8 of the specification
derives signing randomness deterministically from \(K\) and \(\mu\).  The formal
wrapper uses \ec{drandom\_coins} and a typed sampler ROM query.  The security
proof is therefore a theorem about this modeled interface, and the dissertation
records the remaining obligation to relate it to the deterministic specification
path.  The main text now states this relation as the explicit modeling
assumption in \Cref{assm:signing-randomness-modeling-bridge}; this appendix
should be read as code-level evidence for the EasyCrypt side of that bridge, not
as a machine-checked proof of the bridge itself.

\section{Parameters as mathematical constants}
\label{app:case-params}

\begin{lstlisting}[language=EasyCrypt,caption={Representative parameter declarations in \ec{HAETAE\_Params.ec}.},label={lst:haetae-params-snippet}]
op seedbytes : int = 32.
op crhbytes : int = 64.
op n : int = 256.
op q : int = 64513.
op dq : int = 2 * q.

type mode = [ Mode2 | Mode3 | Mode5 ].

op mode_k (md : mode) : int =
  with md = Mode2 => 2
  with md = Mode3 => 3
  with md = Mode5 => 4.

op mode_l (md : mode) : int =
  with md = Mode2 => 4
  with md = Mode3 => 6
  with md = Mode5 => 7.

op mode_tau (md : mode) : int =
  with md = Mode2 => 58
  with md = Mode3 => 80
  with md = Mode5 => 128.
\end{lstlisting}

Mathematically, \ec{mode} indexes three parameterized schemes
\(\haetae_2\), \(\haetae_3\), and \(\haetae_5\).  The functions \ec{mode\_k},
\ec{mode\_l}, and \ec{mode\_tau} are not runtime computations in the proof
argument; they are mode-indexed constants.  Lemmas over \ec{haetae\_mode}
therefore mean theorems for the selected parameter instance.

\section{The counted-ROM bound}
\label{app:case-counted-rom-bound}

\begin{lstlisting}[language=EasyCrypt,caption={Final counted-ROM bound expression.},label={lst:case-counted-rom-bound}]
op counted_rom_programming_loss_term =
  counted_rom_collision_term +
  counted_rom_prequery_term +
  counted_rom_min_entropy_term.

op haetae_euf_counted_rom_bound =
  mlwe_hardness_term +
  bimodal_selftarget_msis_hardness_term +
  counted_rom_programming_loss_term.
\end{lstlisting}

The mathematical conclusion of the final fixed-ledger composition lemmas is
\[
  \mathcal{H}_{\mathsf{hop}}
  \wedge
  \mathcal{H}_{\mathsf{hard}}
  \Longrightarrow
  \Adv_{\mathsf{ECModel},\rom}^{\eufcma}(A)
  \le
  L_{\mathsf{MLWE}}
  +L_{\mathsf{ModSIS}}
  +L_{\mathsf{bimodal}\to\mathsf{MSIS}}
  +L_{\mathsf{coll}}
  +L_{\mathsf{pre}}
  +L_{\mathsf{me}}.
\]
The EasyCrypt operation \ec{bimodal\_selftarget\_msis\_hardness\_term} expands
to \ec{module\_sis\_hardness\_term + bimodal\_to\_msis\_reduction\_loss\_term}.
Thus the final expression is a concrete algebraic object, not an informal bound
name.

\section{Lazy-ROM freshness}
\label{app:case-lazy-rom-freshness}

\begin{lstlisting}[language=EasyCrypt,caption={Concrete freshness lemma for a sampler ROM point.},label={lst:case-fresh-le}]
lemma concrete_lazy_rom_sampler_expand_query_fresh_le
    seed_coins (p : random_coins -> bool) &m :
  sampler_expand_query seed_coins \notin HAETAE_RO.FRO.m{m} =>
  Pr[HAETAE_RO.FRO.get(sampler_expand_query seed_coins) @ &m :
       p (ro_signing_coins res)] <=
  mu drandom_coins p.
\end{lstlisting}

This lemma says that if the sampler query is fresh for the lazy ROM table, then
asking the random oracle at that point produces signing coins distributed no
worse than the ideal \ec{drandom\_coins} measure for every predicate \(p\).  In
measure-theoretic notation, for a fresh point \(q\),
\[
  \Pr[p(\mathsf{ro\_signing\_coins}(H(q)))]
  \le
  \mu_{D_{coins}}(p).
\]
The proof is the local semantic reason that the simulator may treat an
unqueried sampler point as fresh randomness.

\section{Active signing and clean lifting}
\label{app:case-active-signing-clean-lifting}

\begin{lstlisting}[language=EasyCrypt,caption={Clean sampler-ROM preservation across one active signing call.},label={lst:case-o-sign-clean}]
lemma concrete_budgeted_ro_signing_attempt_o_sign_clean_sampler_fresh :
  hoare[ROMInternalTranscriptBudgetedPaperSimAsNMA
          (A, ROSigningAttemptPaperSimSampler(HAETAE_RO.FRO),
           HAETAE_RO.FRO).O.sign :
    sampler_rom_covered
      HAETAE_RO.FRO.m
      ROMInternalTranscriptBudgetedPaperSimAsNMA.adversary_hash_queries
      ROMInternalTranscriptBudgetedPaperSimAsNMA.sampler_expand_queries /\
    ! ROMInternalTranscriptBudgetedPaperSimAsNMA.sampler_bad_prequery ==>
    sampler_rom_covered
      HAETAE_RO.FRO.m
      ROMInternalTranscriptBudgetedPaperSimAsNMA.adversary_hash_queries
      ROMInternalTranscriptBudgetedPaperSimAsNMA.sampler_expand_queries].
\end{lstlisting}

This displayed header is the local theorem surface needed for the mathematical
reading:
\[
  \{C\}\;O_{\mathsf{attempt}}.\mathsf{sign}\;\{C\},
\]
where \(C\) contains sampler-ROM coverage and the absence of
\ec{sampler\_bad\_prequery}.  This Hoare judgment establishes that the clean
condition needed for the one-call surface lemma remains available after a
signing query.  By itself it does not prove that the final NMA theorem path is
non-vacuous; the later composition still has to consume this local surface
instead of relying on coarse fallback probability bounds.

\section{Hoare and equivalence judgments as proof objects}
\label{app:case-hoare-equiv}

A Hoare judgment
\[
  \ec{hoare}[c:P==>Q]
\]
is a partial-correctness statement: every terminating execution of \(c\) from a
\(P\)-state ends in a \(Q\)-state.  When \(c\) is additionally proved lossless,
the same support-inclusion fact can be read as probability-one postcondition
preservation.  An equivalence
judgment
\[
  \ec{equiv}[c_1 \sim c_2:P==>Q]
\]
is a coupling statement between two probabilistic programs.  In the HAETAE proof,
Hoare judgments maintain invariants such as budget bounds and clean ROM state,
while equivalence judgments justify exact game hops.  Probability lemmas then
turn bad-event bounds and reduction assumptions into the final advantage
inequality.

\section{Why these cases are representative}
\label{app:case-representative}

The full proof contains many declarations, but the main mathematical burden is
carried by a small number of object classes: the signature-game interface, the
HAETAE scheme wrapper, mode-indexed parameters, lazy-ROM freshness lemmas,
active-signing Hoare invariants, sampler/exactness lifting lemmas, and the final
bound operations.  The archived exhaustive catalog can be used to audit names;
the curated cases above explain why the names matter mathematically.
